\section{\ENG{Top-Level Runtime Architecture}}
\label{sec:top_level_runtime_architecture}

Στο υψηλότερο επίπεδο αφαίρεσης, ο \ENG{runtime core} του \ENG{engine} οργανώνεται γύρω από τρεις βασικές μεταβάσεις: δημιουργία \ENG{launch}, ανάκτηση \ENG{launch context} και λήψη \ENG{runtime commits}. Η πρώτη μεταβαίνει από το \ENG{LMS} στη συνεδρία παρακολούθησης. Η δεύτερη τροφοδοτεί τον \ENG{player} με τα δεδομένα που απαιτούνται για να συνθέσει το \ENG{launch shell}. Η τρίτη επιτρέπει την προοδευτική ενημέρωση της κατάστασης εκτέλεσης μέχρι τον τελικό τερματισμό της προσπάθειας.

Η δημιουργία \ENG{launch} υλοποιείται από τον \texttt{\ENG{CreateLaunchCommandHandler}}. Εκεί ελέγχεται η ύπαρξη ενεργής εγγραφής του χρήστη στο μάθημα, εντοπίζεται πιθανή ήδη ενεργή προσπάθεια και, ανάλογα με την παράμετρο \texttt{\ENG{forceNewAttempt}}, είτε απορρίπτεται η νέα αίτηση είτε κλείνεται η παλαιά προσπάθεια και δημιουργείται νέα. Στη συνέχεια δημιουργούνται τόσο η εγγραφή \ENG{attempt} όσο και η εγγραφή \ENG{launch}, εκδίδεται \ENG{launch token} μέσω του \texttt{\ENG{JwtTokenService}} και αρχικοποιείται η \ENG{LaunchRuntimeState} στο \texttt{\ENG{RuntimeStateStore}}. Με τον τρόπο αυτό, το \ENG{launch} είναι πλήρως ορισμένο \ENG{domain} γεγονός που συνδέει \ENG{user}, \ENG{course}, \ENG{attempt} και κατάσταση εκτέλεσης.

Η ανάκτηση του \ENG{launch context} υλοποιείται από τον \texttt{\ENG{LaunchOrchestrator}}. Η μέθοδος \texttt{\ENG{getLaunchContext()}} επαληθεύει ότι το \ENG{launch} είναι ακόμη ενεργό, ότι το \ENG{launch token} ταιριάζει με το ζητούμενο \ENG{launchId} και ότι δεν έχει λήξει χρονικά. Κατόπιν συνθέτει το \texttt{\ENG{LaunchContextDto}} με στοιχεία χρήστη, μαθήματος, \ENG{player} ρύθμισης, \ENG{presigned ZIP URL} από το \ENG{object storage} και την τρέχουσα κανονικοποιημένη κατάσταση της συνεδρίας. Η ύπαρξη του \texttt{\ENG{RuntimeView(initialNormalizedState)}} στο \ENG{DTO} είναι σημαντική, διότι δείχνει ότι ο \ENG{engine} είναι ικανός να εκθέσει την ήδη γνωστή πρόοδο, ακόμη κι αν ο \ENG{player} δεν την ζητά ακόμη στο \ENG{browser-side API state}.

Η λήψη \ENG{commit requests} υλοποιείται από τον \texttt{\ENG{RuntimeOrchestrator}}. Σε αυτή τη φάση ο πυρήνας επαληθεύει την κατάσταση του \ENG{launch}, ελέγχει το \ENG{launch-token scope}, εφαρμόζει \ENG{rate limiting}, ανακτά ή αρχικοποιεί \ENG{LaunchRuntimeState}, απορρίπτει παρωχημένες ακολουθίες \ENG{sequence numbers} όταν το επιβάλλει η ρύθμιση \texttt{\ENG{strictCommitSequence}}, και μόνο τότε εκχωρεί την ερμηνεία του \ENG{payload} σε \ENG{adapter} που αντιστοιχεί στο συγκεκριμένο πρότυπο. Η ροή αυτή καθιστά σαφές ότι ο \ENG{engine} μετασχηματίζει τα δεδομένα πρώτα σε εσωτερική κατάσταση συνεδρίας.

Το Σχήμα~\ref{fig:top_level_runtime} συνοψίζει το τοπολογικό αυτό μοτίβο. Ο Πίνακας~\ref{tab:engine_endpoint_responsibility} συσχετίζει τις βασικές \ENG{REST} επιφάνειες με τους αντίστοιχους \ENG{orchestrators} και τα όρια μόνιμης αποθήκευσης, ώστε να γίνει σαφές σε ποιο σημείο του συστήματος λαμβάνεται κάθε κρίσιμη απόφαση.

\begin{figure}[H]
  \centering
\begin{chapterfigurebox}
  \begin{tikzpicture}[
      font=\footnotesize,
      node distance=1.8cm and 2.4cm,
      box/.style={draw, rounded corners, fill=gray!5, align=center, text width=3.2cm, minimum height=1.1cm},
      emph/.style={draw, rounded corners, fill=gray!12, align=center, text width=3.5cm, minimum height=1.1cm, font=\footnotesize\bfseries},
      store/.style={draw, cylinder, shape border rotate=90, aspect=0.28, minimum height=1.2cm, minimum width=3.2cm, align=center, fill=gray!5},
      arrow/.style={-Latex, thick}
    ]
    \node[emph] (player) {\ENG{SCORM Player}};
    \node[box, below left=1.8cm and 1.8cm of player] (launchc) {\texttt{\ENG{LaunchController}}};
    \node[box, below right=1.8cm and 1.8cm of player] (runtimec) {\texttt{\ENG{RuntimeController}}};
    \node[emph, below=of launchc] (launcho) {\texttt{\ENG{LaunchOrchestrator}}};
    \node[emph, below=of runtimec] (runtimeo) {\texttt{\ENG{RuntimeOrchestrator}}};
    \node[box, below=of runtimeo] (registry) {\texttt{\ENG{StandardAdapterRegistry}}};

    \node[store, below=2.0cm of launcho] (postgres) {\ENG{Postgres}\\\ENG{attempts / launches}};
    \node[store, right=of postgres] (redis) {\ENG{Redis}\\κατάσταση εκτέλεσης};
    \node[store, right=of redis] (minio) {\ENG{MinIO}\\\ENG{ZIP content}};

    \draw[arrow] (player) -- (launchc);
    \draw[arrow] (player) -- (runtimec);
    \draw[arrow] (launchc) -- (launcho);
    \draw[arrow] (runtimec) -- (runtimeo);
    \draw[arrow] (runtimeo) -- (registry);
    \draw[arrow] (launcho) -- (postgres);
    \draw[arrow] (launcho) -- (redis);
    \draw[arrow] (launcho) -- (minio);
    \draw[arrow] (runtimeo) -- (redis);
    \draw[arrow] (runtimeo) |- (postgres);
  \end{tikzpicture}
\end{chapterfigurebox}
  \caption{Υψηλού επιπέδου αρχιτεκτονική του \ENG{runtime} πυρήνα, από τον \ENG{player} έως τα κύρια όρια μόνιμης αποθήκευσης.}
  \label{fig:top_level_runtime}
\end{figure}


\begin{table}[H]
\centering
\caption{Κύριες \ENG{engine} διεπαφές και αντιστοίχισή τους με τους εσωτερικούς μηχανισμούς επεξεργασίας}
\label{tab:engine_endpoint_responsibility}
\begingroup
\footnotesize
\setlength{\tabcolsep}{3pt}
\renewcommand{\arraystretch}{1.20}
\begin{adjustbox}{center,max width=\textwidth}
\begin{tabular}{|p{3.8cm}|p{3.2cm}|p{4.1cm}|p{3.6cm}|}
\hline
\ENG{Endpoint group} & Κύριος \ENG{controller} & Κεντρικός \ENG{handler} / \ENG{orchestrator} & Κύριο όριο μόνιμης αποθήκευσης \\
\hline
\texttt{\ENG{/api/v1/courses/import}} & \texttt{\ENG{CourseController}} & \texttt{\ENG{ImportCourseCommandHandler}}, \texttt{\ENG{ManifestParser}} & \ENG{MinIO} για \ENG{ZIP}, \ENG{Postgres} για \texttt{\ENG{courses}} \\
\hline
\texttt{\ENG{/api/v1/users}}, \texttt{\ENG{/api/v1/enrollments}} & \texttt{\ENG{UserController}}, \texttt{\ENG{EnrollmentController}} & Υπηρεσίες ανάγνωσης/καταχώρισης των αντίστοιχων \ENG{domain objects} & \ENG{Postgres} \\
\hline
\texttt{\ENG{/api/v1/launches}} & \texttt{\ENG{LaunchController}} & \texttt{\ENG{CreateLaunchCommandHandler}} & \ENG{Postgres} για \texttt{\ENG{attempts}}/\texttt{\ENG{launches}}, \ENG{Redis} για αρχική κατάσταση \\
\hline
\texttt{\ENG{/api/v1/launches/\{id\}}} & \texttt{\ENG{LaunchController}} & \texttt{\ENG{LaunchOrchestrator}} & \ENG{Postgres} για αναγνώσεις επιχειρησιακών δεδομένων, \ENG{Redis} για τρέχουσα κατάσταση, \ENG{MinIO} \ENG{presigned URL} \\
\hline
\texttt{\ENG{/api/v1/runtime/launches/\{id\}/commit}} & \texttt{\ENG{RuntimeController}} & \texttt{\ENG{RuntimeOrchestrator}}, \texttt{\ENG{StandardAdapterRegistry}} & \ENG{Redis} για ενεργή κατάσταση, \texttt{\ENG{launches.lastSeenAt}} στο \ENG{Postgres} \\
\hline
\texttt{\ENG{/api/v1/launches/\{id\}/terminate}} & \texttt{\ENG{LaunchController}} & \texttt{\ENG{TerminateLaunchCommandHandler}}, \texttt{\ENG{RuntimeFlushService}} & \ENG{Postgres} για κλείσιμο προσπάθειας/στιγμιότυπο, διαγραφή από \ENG{Redis} \\
\hline
\texttt{\ENG{/api/v1/attempts/*}}, \texttt{\ENG{/api/v1/courses/*/stats}} & \texttt{\ENG{AttemptController}}, \texttt{\ENG{ReportingController}} & \ENG{Read/report services} & \ENG{Postgres} \\
\hline
\end{tabular}
\end{adjustbox}
\endgroup

\end{table}
